\documentclass[14pt, oneside]{article} 
\usepackage{geometry}
\geometry{letterpaper} 
\usepackage{graphicx}
	
\usepackage{amssymb}
\usepackage{amsmath}
\usepackage{parskip}
\usepackage{color}
\usepackage{hyperref}

\graphicspath{{/Users/telliott/Github-math/figures/}}
% \begin{center} \includegraphics [scale=0.4] {gauss3.png} \end{center}

\title{Inversive Geometry}
\date{}

\begin{document}
\maketitle
\Large

%[my-super-duper-separator]

\emph{Inversion} in a circle maps a point in the plane to (typically) a different point in the plane.  The mapping is continuous, so if points lie close together before inversion they are close together afterwards.  A locus of points is transformed into another locus, e.g. a circle is in most cases transformed into another circle.

To carry out inversion, one starts by choosing a reference circle.  We will refer to $\omega$, the \emph{inversion circle}, with its \emph{inversion center} at $O$ having radius $k$.  We show $\omega$ as a dotted circle.

Let $P$ be an arbitrary point, distinct from $O$.  The inverse of $P$ is called $P'$.  $P'$ lies on the ray through $OP$ according to the following rule:
\begin{center} \includegraphics [scale=0.40] {inversion1.png} \end{center}
\[ OP \cdot OP' = k^2 \ \ \ \ \ \ \ \  \frac{OP}{k} = \frac{k}{OP'} \]

We can write $ I_{\omega} (P) = P'$.  

Rather than the ray, we will just draw a line segment, with definite endpoints, and also as usual just call it line segment a line, for brevity.

A locus of transformed points is referred to as the image of the input locus under inversion.  In some cases the image of a point is the same as the original.  Points lying on $\omega$ invert into themselves, because then $OP = k = OP'$.

You can see in the example above that $OP$ is about one-half of $k$, so then the length of $OP'$ is twice $k$.  The image of a point close to $O$ may be very far away.

$O$ is a special point in more than one way.  As $P$ approaches $O$, the inverse $P'$ moves farther and farther away.

Here's an overview, with the \emph{inversion circle} represented as a dotted circle.  The blue circle maps to the one in cyan and the red circle maps to the one in tan.  The converse is also true.
\begin{center} \includegraphics [scale=0.55] {inversion16.png} \end{center}  

Points outside the inversion circle get mapped to points inside, and vice-versa.  A line through the center maps to itself.

A locus through the center of the inversion circle (magenta circle) is a special case --- it maps to the purple line. The black circle is another special case, it maps to a circle that looks just like original, although all of the points have been swapped (except the two lying on the inversion circle).

The complete theory invents a "point at infinity" which is the inverse of $O$.  We won't worry about that.  We'll just think of the set of possible inputs to $I$ as the "punctured plane", with $O$ missing.

When we say something like ``a circle through $O$'', we are not worrying about the point $O$ itself.
\begin{center} \includegraphics [scale=0.30] {inversion1a.png} \end{center}

We call points that are invariant under inversion, fixed.  Let $R$ be a point on $\omega$, so the length of $OR$ is $k$.  Then we have that
\[ I_{\omega} (R) = R' = R \]

Much of our work with inversions involves similar triangles.  Coxeter introduces it this way:
\begin{center} \includegraphics [scale=0.30] {Coxeter_5_3A.png} \end{center}
Draw the chord $TU \perp OP$.  Then $P'$ is the point where the tangents from $T$ and $U$ meet.  From similar triangles $\triangle OTP' \sim \triangle OPT$, so we have:
\[ \frac{OP'}{k} = \frac{k}{OP} \]

An analogous transformation in three dimensions uses a sphere, in space.

\subsection*{inversion is an involution}

If $P'$ is the inverse of $P$, then the inverse of $P'$ is $P$.

\emph{Proof}.  Let $P''$ be the inverse of $P'$, with $OP' \cdot OP'' = k^2$.  But $k^2 = OP \cdot OP'$.  It follows that $OP'' = OP$. $\square$

If $P$ lies inside $\omega$, then $P'$ lies outside, and vice-versa.  There is a famous essay which analogizes how one would hunt a lion in the desert by various mathematical or physical techniques.  Here is the one using inversion:

\begin{quote}
We place a \emph{spherical} cage in the desert, enter it and lock it. We perform an inversion with respect to the cage. The lion is then inside the cage, and we are outside.  \emph{H. Petard}
\end{quote}

\subsection*{inversion of a line through $O$}

We now consider the results of inverting either a line or a circle.  There are two cases for each, depending on whether the line or circle goes through, $O$, the center of $\omega$.  They are all related, as we will see.

\textbf{The inverse of a line through $O$ is the same line}.

\begin{center} \includegraphics [scale=0.25] {inversion2.png} \end{center}
Let $P$ and $Q$ be any two points on a line through $O$.  The inverse points $P'$ and $Q'$ lie on the same line.

\subsection*{inversion of a line not through $O$}

\begin{center} \includegraphics [scale=0.3] {inversion3a.png} \end{center}

\textbf{The inverse of a line not through $O$ is a circle through $O$.}

This is the first result that is a little surprising.

\emph{Proof}.

Let $L$ be a line, not passing through $O$ and external to $\omega$, and let $P$ and $A$ be two points on the line.  Pick $A$ such that $OA'A \perp PA$.

We have that 
\[ OP \cdot OP' = k^2 = OA \cdot OA' \]
\[ \frac{OP}{OA} = \frac{OP'}{OA'} \]

It follows that $\triangle PAO \sim \triangle A'P'O$ and in particular, $\angle OP'A'$ is a right angle because $\angle OAP$ is a right angle.

But this is true for \emph{every} point on $L$ other than $A$.

By the inverse of Thales' circle theorem, it follows that $A'$ and $P'$ lie on a circle.  This circle passes through $O$ and in fact, has diameter $OA'$.

$\square$

The diameter of this circle is perpendicular to the line $L$, so $L$ is parallel to the tangent to the circle at $O$ (not shown).

\begin{center} \includegraphics [scale=0.3] {inversion4.png} \end{center}

The same proof still works for the case where $L$ passes through $\omega$, or even when it is tangent (below).  This is because, in both cases, $\triangle OAP \sim \triangle OP'A'$.

\begin{center} \includegraphics [scale=0.3] {inversion4b.png} \end{center}

\subsection*{note on similar triangles}

For any points $A$ and $P$, we can form similar triangles by using the fundamental equations
\[ OP \cdot OP' = k^2 = OA \cdot OA' \]
\[ \frac{OP}{OA} = \frac{OP'}{OA'} \]

\begin{center} \includegraphics [scale=0.38] {inversion3b.png} \end{center}

This works even when one of $A$ or $P$ lies inside, and one outside $\omega$.

\subsection*{inversion of a circle through $O$}

Let $K$ be a circle, passing through $O$.  

\begin{center} \includegraphics [scale=0.3] {inversion3a.png} \end{center}

\textbf{The inverse of a circle through $O$ is a line not through $O$.}

This follows directly from the previous theorem, since the inverse of the inverse of any point is the same point.

As we saw from the previous result, it does not matter whether the circle is entirely internal to $\omega$, is internally tangent, or extends outside.

\subsection*{Ptolemy's Theorem}

With only what we have established so far, we can prove an important theorem.

Let $A$, $B$ and $C$ lie on a circle through $O$.  Then, as we showed above, the inverse of this circle is a line not through $O$ and here, entirely external to the inversion circle.

Ptolemy says that 
\[ AB \cdot OC + BC \cdot OA = AC \cdot OB \]

\emph{Proof}.

\begin{center} \includegraphics [scale=0.35] {inversion9.png} \end{center}

$A'B' + B'C' = A'C'$.  We will leverage this identity to find the result.

We have three pairs of similar triangles.  To help stay organized, we write out the ratio boxes.
\begin{center} \includegraphics [scale=0.18] {ratios11.png} \end{center}

The symmetry is pretty clear.  We have that
\[ \frac{A'B'}{AB} = \frac{OA'}{OB} \Rightarrow A'B' =  \frac{AB \cdot OA'}{OB} \]

Also
\[ \frac{B'C'}{BC} = \frac{OC'}{OB} \Rightarrow B'C' = \frac{BC \cdot OC'}{OB} \]

Then
\[ \frac{A'C'}{AC} = \frac{OC'}{OA} \Rightarrow A'C' = \frac{AC \cdot OC'}{OA} \]

Form the sum and clear the denominators:
\[ AB \cdot OA' \cdot OA + BC \cdot OC' \cdot OA = AC \cdot OC' \cdot OB \]

Divide by $OC'$
\[ AB \cdot \frac{OA'}{OC'} \cdot OA + BC \cdot OA = AC \cdot OB \]

Can you find the appropriate entry for the last step?  I obtain
\[ AB \cdot OC + BC \cdot OA = AC \cdot OB \]

$\square$

This is Ptolemy's theorem.

We pause to reflect on one of the most useful and important theorems in all of geometry.

We now consider circles not through $O$.  This is the point where the theory gets a little more sophisticated.

\begin{center} \includegraphics [scale=0.55] {inversion15.png} \end{center}

\subsection*{inversion of a circle not through $O$}

\textbf{The inverse of a circle not through $O$ is another circle not through $O$.}

We can again draw different diagrams depending on where the circles lie with respect to $\omega$.  We need to make sure that the proof or proofs work for all of them.

The cases we've already seen include a circle through the inversion center, which inverts into a line extending infinitely out (magenta and purple).

Consider first the case where $K$ is entirely external to $\omega$.

\emph{Proof}.

Let $K$ be a circle, not passing through $O$, and not even intersecting with $\omega$.  Draw the diameter $AB$ of $K$, such that the extension passes through $O$.

\begin{center} \includegraphics [scale=0.35] {inversion5.png} \end{center}

Let $C$ be an otherwise arbitrary point on $K$, not on $AB$.  Find the inverses of $A$, $B$ and $C$.  Draw $\triangle A'B'C'$.

As before
\[ \frac{OA'}{OC'} = \frac{OA}{OC} \]
so $\triangle OB'C' \sim \triangle OCB$ and $\triangle OA'C' \sim \triangle OCA$.

From similar triangles, we have that the angles marked $\alpha$ and $\beta$ are equal.

By Thales' circle theorem, $\angle ACB$ is a right angle.  So from $\triangle BOC$:
\[ \angle BOC + \alpha + \beta + \angle ACB = 180 \]

But from $\triangle B'OC'$ (with $\angle B'OC' = \angle BOC$):
\[ \angle B'OC'  + \alpha + \beta + \angle AC'B' = 180 \]

It follows that $\angle A'C'B' = \angle ACB$, so $\angle A'C'B'$  is also a right angle, and thus it lies on the circle whose diameter is $A'B'$.  This is true for every point $C$ on circle $K$, other than $A$ and $B$.  

$\square$

Second case: let the circle $K$ intersect $\omega$, with the inversion center $O$ inside $K$.

Draw the diameter $AB$ of circle $K$ which goes through $O$.  Let $C$ be any other point on circle $K$.

\begin{center} \includegraphics [scale=0.36] {inversion6.png} \end{center}

\emph{Proof}.

We have $\triangle OAC \sim \triangle OC'A'$ and $\triangle OBC \sim \triangle OC'B'$.  From similar triangles, the angles marked as $\alpha$ are equal, as are those marked as $\beta$.

$\angle ACB$ is right, but $\alpha + \beta$ is supplementary to $\angle ACB$ to $\alpha + \beta$ so the sum is $90$.

It follows that $\angle A'C'B' = \alpha + \beta$ is also right.  Therefore $C'$ lies on a circle whose diameter is $A'B'$.  

$\square$.

Third case: the circle $K$ intersects $\omega$ but the inversion center $O$ is external to $K$.

\begin{center} \includegraphics [scale=0.36] {inversion7.png} \end{center}

Draw the diameter $AB$ of circle $K$ which, extended, goes through $O$.  Let $C$ be any other point on circle $K$.

\emph{Proof}.

The angles marked $\alpha$ and $\beta$ are equal (in pairs) by similar triangles.

$\angle ACB$ is right.  So $\angle BOC + \alpha + \beta$ is equal to a right angle.

But $\angle BOC + \alpha + \beta + \angle A'B'C'$ is equal to two right angles.

Thus, $\angle A'C'B'$ is also right.

$\square$

The fourth and last case is where the circle $K$ and $\omega$ are concentric.  But then the inverse of $K$ is also a concentric circle.


\end{document}